% Рисунок с доски: обобщённый материал Кельвина — Фойгта — последовательная
% цепочка из N звеньев «пружина ∥ демпфер».
\begin{figure}[htbp]
\centering
\begin{tikzpicture}[>={Latex[length=4pt]}, scale=0.95, transform shape]
  \foreach \k/\lab in {0/1, 3.2/2, 6.4/3} {
    \begin{scope}[xshift=\k cm]
      \pic at (0,0) {spring};
      \node[font=\footnotesize, anchor=north] at (0,-0.45) {$\mu_{\lab}$};
      \pic at (0,1.2) {damper};
      \node[font=\footnotesize, anchor=south] at (0,1.55) {$\eta_{\lab}$};
      \draw[thick] (-0.9,0) -- (-1.2,0) -- (-1.2,1.2) -- (-0.9,1.2);
      \draw[thick] (0.9,0) -- (1.2,0) -- (1.2,1.2) -- (0.9,1.2);
      \node[font=\footnotesize, anchor=west] at (1.3,0.6)
        {$\tau_{\lab} = \frac{\eta_{\lab}}{\mu_{\lab}}$};
    \end{scope}
  }
  \draw[thick] (-1.2,0.6) -- (-1.8,0.6);
  \draw[thick] (2.0,0.6) -- (2.3,0.6);
  \draw[thick] (5.2,0.6) -- (5.5,0.6);
  \node at (8.6,0.6) {$\cdots$};
  \begin{scope}[xshift=10.6cm]
    \pic at (0,0) {spring};
    \node[font=\footnotesize, anchor=north] at (0,-0.45) {$\mu_N$};
    \pic at (0,1.2) {damper};
    \node[font=\footnotesize, anchor=south] at (0,1.55) {$\eta_N$};
    \draw[thick] (-0.9,0) -- (-1.2,0) -- (-1.2,1.2) -- (-0.9,1.2);
    \draw[thick] (0.9,0) -- (1.2,0) -- (1.2,1.2) -- (0.9,1.2);
    \draw[thick] (1.2,0.6) -- (1.8,0.6);
  \end{scope}
  \draw[thick] (9.0,0.6) -- (9.4,0.6);
\end{tikzpicture}
\caption{Обобщённый материал Кельвина–Фойгта: $N$ звеньев соединены
последовательно, поэтому $\dev = \sum_k \dev_k$ при общем $\Sdev$, а функции
ползучести складываются.}
\label{fig:kv-generalized}
\end{figure}
